\begin{abstract}
\noindent
The key--value (KV) cache is the dominant memory cost of long-context transformer inference: it
grows linearly with sequence length and, on a single commodity device, becomes the binding
constraint long before compute does. We present \diffkv{}, a KV-cache inference runtime that
attacks this cost structurally rather than by eviction or quantization alone. \diffkv{} partitions
the cache into fixed-size blocks and represents each block three ways at once: an \emph{anchor}
token, a low-rank \emph{differential} term (a rank-$16$ joint randomized SVD of the per-token
deviations from the anchor, shared across keys and values), and a small set of \emph{exact residual
tokens}---the handful of tokens whose low-rank reconstruction error is largest, kept verbatim in
fp16. The low-rank term captures a block's shared structure cheaply; the residuals restore the
low-energy outliers (a rare entity, a buried passcode) that pure low-rank compression discards.
A sliding dense \emph{recency window} keeps the most recent tokens uncompressed, and the whole
store is a pre-allocated, fixed-size pool whose footprint is \emph{bounded and context-independent}
($\approx\!0.68$\,GB across $28$ layers), while a dense full-KV cache grows without bound. At
decode, a single fused, compiled kernel routes the query to the most relevant blocks by an
\emph{exact} residual-key score, evaluates those blocks' attention \emph{directly in low-rank
space}---never reconstructing dense keys or values---attends their exact residuals together with
the recency window, and merges the two branches with a numerically stable log-sum-exp. An adaptive
policy uses exact full-KV decode at short context and switches to the compressed kernel only once
the cache is large enough to be the binding cost.

We implement \diffkv{} on Apple's MLX framework and evaluate it on Qwen2.5-1.5B-Instruct (int4) on
an $8.6$\,GB Apple~M3, with a needle-in-a-haystack protocol from $4$k to $64$k tokens. We report
results honestly and separate the compressed sparse decode---the runtime's actual
contribution---from an exact full-KV decode used as an upper-bound ablation on the same store.
\diffkv{}'s compressed KV state is $\approx\!2.85\times$ smaller than the dense cache at equal
precision and remains bounded as context grows, letting the model process a $64$k prompt on a
device where the dense baseline exhausts memory by $32$k---and the residual mechanism makes the
compressed decode recover the buried passcode \emph{exactly} at every tested context. These
benefits come at measured costs we quantify rather than hide: the compressed decode is several
times slower than dense decode (an implementation gap in the kernel dispatch, localized by the
exact ablation), and the per-block SVD makes prefill the slowest phase. We characterize the central
residual-budget trade-off---more exact residuals buy recall at the cost of compression---and design
every figure and table so a future CUDA/Triton fused-decode implementation can be inserted without
restructuring the evaluation. A full reproducibility appendix accompanies the report; all reported
numbers are measured on the same machine, and no value is averaged across conflicting runs or
fabricated.
\end{abstract}
